\documentclass[10pt, aspectratio=169]{beamer}
% Preámbulo modular para presentaciones Beamer (Modelación Estadística)
% Diseño y paleta institucional del Tecnológico de Monterrey

\documentclass[aspectratio=169,xcolor={dvipsnames,table}]{beamer}

\usepackage[utf8]{inputenc}
\usepackage[spanish,mexico]{babel}
\usepackage{amsmath,amssymb,amsfonts,mathtools}
\usepackage{graphicx}
\usepackage{listings}
\usepackage{booktabs}
\usepackage{tikz}
\usetikzlibrary{matrix,arrows,decorations.pathmorphing,shapes.geometric,positioning}

% Paleta Institucional Tecnológico de Monterrey
\definecolor{TecAzul}{HTML}{0039A6}       % Azul TEC: color primario institucional
\definecolor{TecAzulOscuro}{HTML}{1A2E51} % Azul marino oscuro
\definecolor{TecRojo}{HTML}{EC2661}       % Rojo de acento (paleta miTec)
\definecolor{TecGris}{HTML}{646464}       % Gris neutro para comentarios y subtítulos
\definecolor{TecFondoBloque}{HTML}{F4F6F9}% Fondo suave para bloques

% Configuración de Tema Beamer
\usetheme{Boadilla}
\usecolortheme[named=TecAzulOscuro]{structure}
\setbeamercolor{alerted text}{fg=TecRojo}
\setbeamercolor{palette primary}{bg=TecAzulOscuro,fg=white}
\setbeamercolor{palette secondary}{bg=TecAzul,fg=white}
\setbeamercolor{palette tertiary}{bg=TecAzulOscuro!80!black,fg=white}
\setbeamercolor{title}{fg=TecAzulOscuro,bg=TecFondoBloque}
\setbeamercolor{frametitle}{fg=TecAzulOscuro,bg=TecFondoBloque}

% Bloques personalizados elegantes
\setbeamercolor{block title}{bg=TecAzulOscuro,fg=white}
\setbeamercolor{block body}{bg=TecFondoBloque,fg=black}
\setbeamercolor{block title alerted}{bg=TecRojo,fg=white}
\setbeamercolor{block body alerted}{bg=TecRojo!8,fg=black}
\setbeamercolor{block title example}{bg=TecAzul,fg=white}
\setbeamercolor{block body example}{bg=TecAzul!8,fg=black}

% Redondear bordes y añadir sombra sutil
\setbeamertemplate{blocks}[rounded][shadow=true]
\setbeamertemplate{navigation symbols}{}

% Pie de página limpio con número de diapositiva
\setbeamertemplate{footline}{
  \leavevmode%
  \hbox{%
  \begin{beamercolorbox}[wd=.7\paperwidth,ht=2.25ex,dp=1ex,left]{author in head/foot}%
    \hspace*{2ex}\insertshortauthor~~(\insertshortinstitute) \hfill \insertshorttitle
  \end{beamercolorbox}%
  \begin{beamercolorbox}[wd=.3\paperwidth,ht=2.25ex,dp=1ex,right]{date in head/foot}%
    \insertshortdate{}\hspace*{2em}
    \insertframenumber{} / \inserttotalframenumber\hspace*{2ex}
  \end{beamercolorbox}}%
  \vskip0pt%
}

% Configuración de Listings de Python para visualización pedagógica de código
\lstset{
    language=Python,
    basicstyle=\ttfamily\footnotesize,
    keywordstyle=\color{TecAzulOscuro}\bfseries,
    stringstyle=\color{TecRojo},
    commentstyle=\color{TecGris}\itshape,
    showstringspaces=false,
    breaklines=true,
    frame=lines,
    rulecolor=\color{TecAzulOscuro!30},
    backgroundcolor=\color{TecFondoBloque},
    aboveskip=0.5em,
    belowskip=0.5em,
    literate={á}{{\'a}}1 {é}{{\'e}}1 {í}{{\'i}}1 {ó}{{\'o}}1 {ú}{{\'u}}1
             {Á}{{\'A}}1 {É}{{\'E}}1 {Í}{{\'I}}1 {Ó}{{\'O}}1 {Ú}{{\'U}}1
             {ñ}{{\~n}}1 {Ñ}{{\~N}}1 {¿}{{?`}}1 {¡}{{!`}}1
}

% Metadatos institucionales por defecto
\title[Modelación Estadística]{Modelación Estadística para Ciencia de Datos e Ingeniería}
\author[J. Castillo Colmenares]{Juliho David Castillo Colmenares}
\institute[Tec de Monterrey]{Tecnológico de Monterrey \\ Escuela de Ingeniería y Ciencias}
\date{\today}

% Comandos y macros estadísticas compartidas para las presentaciones Beamer (Metropolis)

% Conjuntos numéricos básicos
\newcommand{\R}{\mathbb{R}}
\newcommand{\N}{\mathbb{N}}
\newcommand{\Q}{\mathbb{Q}}
\newcommand{\Z}{\mathbb{Z}}
\newcommand{\set}[1]{\left\{ #1 \right\}}
\newcommand{\abs}[1]{\left|#1\right|}
\newcommand{\norm}[1]{\left\| #1 \right\|}

% Comandos estadísticos y probabilísticos del curso
\newcommand{\Var}{\operatorname{Var}}
\newcommand{\cov}{\operatorname{Cov}}
\newcommand{\corr}{\rho}
\newcommand{\comb}[2]{\begin{pmatrix} #1 \\ #2 \end{pmatrix}}
\newcommand{\s}{\sigma}
\newcommand{\particion}{\mathfrak{P}}
\newcommand{\card}[1]{\#\left( #1 \right)}
\newcommand{\seq}[3]{\set{#1}_{#2}^{#3}}
\newcommand{\E}{\mathbb{E}}
\newcommand{\Prob}{\mathbb{P}}

% Entornos pedagógicos y cajas en español académico natural
\newenvironment{discusion}{%
  \begin{alertblock}{Pregunta de Discusión}%
}{%
  \end{alertblock}%
}

\newenvironment{reflexion}{%
  \begin{alertblock}{Reflexión para el Aula}%
}{%
  \end{alertblock}%
}

\newenvironment{paradoja}{%
  \begin{alertblock}{Paradoja de Exploración}%
}{%
  \end{alertblock}%
}

\newenvironment{motivacion}{%
  \begin{exampleblock}{Motivación: Ciencia de Datos en la Práctica}%
}{%
  \end{exampleblock}%
}

\newenvironment{retoclase}{%
  \begin{block}{Reto Rápido en Equipo}%
}{%
  \end{block}%
}


\title[TLC Asintótico]{Sección 05.02: Teorema del Límite Central Asintótico}
\subtitle{Demostración vía FGM, Convergencia en Distribución y Teorema de Berry-Esseen}
\author[J. Castillo Colmenares]{Juliho Castillo Colmenares}
\institute{Tecnológico de Monterrey}
\date{\vspace{-1.2cm}}

\begin{document}

% Slide 1: Portada
\begin{frame}[plain]
  \titlepage
\end{frame}

% Slide 2: Hoja de Ruta
\begin{frame}{Hoja de Ruta --- Capítulo 05: Distribuciones de Muestreo}
  \footnotesize
  ¿Dónde nos encontramos en el desarrollo modular del capítulo? \pause
  \begin{itemize}\itemsep=0.08em
    \item \textbf{05.01} Muestreo Aleatorio Simple, Media y Varianza Muestral Insesgada \pause
    \item \textbf{05.02} Teorema del Límite Central Asintótico \emph{(Hoy)} \pause
    \item \textbf{05.03} Distribución Chi-Cuadrada y Varianza Muestral \pause
    \item \textbf{05.04} Distribución $t$ de Student y Muestras Pequeñas \pause
    \item \textbf{05.05} Distribución $F$ de Fisher-Snedecor
  \end{itemize}
  \vspace{-0.05cm}
  \begin{block}{Objetivo de la Sesión}
    Formalizar el Teorema del Límite Central mediante convergencia en distribución, demostrarlo rigurosamente vía la función generadora de momentos, y cuantificar su velocidad de convergencia con el teorema de Berry-Esseen, más allá de la heurística ``$n\ge30$''.
  \end{block}
\end{frame}

% Slide 3: Motivación
\begin{frame}{¿Qué Tan Rápido Converge el TLC?}
  \footnotesize
  \begin{motivacion}
    En la Sección 02.06 vimos el TLC de manera introductoria: la media muestral se aproxima a una Normal cuando $n$ es ``suficientemente grande''. Pero, ¿qué significa exactamente ``suficientemente grande''? ¿Por qué funciona sin importar la distribución original? Hoy formalizamos ambas preguntas.
  \end{motivacion}

  \vspace{0.15cm}
  \pause
  \begin{itemize}\itemsep=0.1cm
    \item \textbf{Convergencia en distribución:} el tipo preciso de límite que aplica al TLC. \pause
    \item \textbf{Demostración vía FGM:} por qué la Normal emerge como límite universal. \pause
    \item \textbf{Berry-Esseen:} una cota explícita del error de la aproximación para $n$ finito.
  \end{itemize}
\end{frame}

% Slide 4: Convergencia en Distribución y TLC
\begin{frame}{Convergencia en Distribución y el Teorema del Límite Central}
  \footnotesize
  \begin{block}{Convergencia en Distribución}
    $Y_n \xrightarrow{d} Y$ si $F_{Y_n}(y) \to F_Y(y)$ en todo punto de continuidad de $F_Y$.
  \end{block}
  \pause

  \vspace{0.1cm}
  \begin{block}{TLC (Lindeberg-Lévy)}
    Si $X_1,X_2,\ldots$ son iid con media $\mu$ y varianza $\sigma^2<\infty$, entonces
    $$Z_n = \dfrac{\bar{X}_n-\mu}{\sigma/\sqrt{n}} \xrightarrow{d} N(0,1) \quad \text{cuando } n\to\infty.$$
  \end{block}

  \vspace{0.1cm}
  \pause
  \begin{alertblock}{Universalidad}
    El resultado \textbf{no depende} de la forma de la distribución original de $X_i$ --- solo requiere media y varianza finitas.
  \end{alertblock}
\end{frame}

% Slide 5: Demostración vía FGM
\begin{frame}{Demostración del TLC vía Función Generadora de Momentos}
  \footnotesize
  Expandimos $\log M_{X_1}(s)$ en serie de Taylor alrededor de $s=0$: \pause

  \vspace{0.1cm}
  \begin{block}{Desarrollo}
    $$M_{Z_n}(t) = e^{-\mu t\sqrt{n}/\sigma}\left[M_{X_1}\!\left(\dfrac{t}{\sigma\sqrt{n}}\right)\right]^n, \qquad \log M_{X_1}(s) = \mu s + \dfrac{\sigma^2 s^2}{2}+O(s^3)$$
  \end{block}

  \vspace{0.1cm}
  \pause
  \begin{alertblock}{Sustituyendo y Tomando el Límite}
    $$\log M_{Z_n}(t) = n\left[\dfrac{\mu t}{\sigma\sqrt{n}}+\dfrac{t^2}{2n}+O(n^{-3/2})\right] - \dfrac{\mu t\sqrt{n}}{\sigma} = \dfrac{t^2}{2}+O(n^{-1/2}) \to \dfrac{t^2}{2}$$
    Por unicidad y continuidad de Lévy, $M_{Z_n}(t)\to e^{t^2/2} \implies Z_n \xrightarrow{d} N(0,1)$.
  \end{alertblock}
\end{frame}

% Slide 6: Teorema de Berry-Esseen
\begin{frame}{Teorema de Berry-Esseen: Tasa de Convergencia}
  \footnotesize
  El TLC garantiza convergencia, pero no indica qué tan rápido. \pause

  \vspace{0.1cm}
  \begin{block}{Cota de Berry-Esseen}
    Con $\rho=E|X_i-\mu|^3<\infty$ y una constante universal $C<0.5$:
    $$\sup_z\left|F_{Z_n}(z)-\Phi(z)\right| \le \dfrac{C\rho}{\sigma^3\sqrt{n}}$$
  \end{block}

  \vspace{0.1cm}
  \pause
  \begin{alertblock}{La Regla ``$n\ge30$'' en Contexto}
    \begin{itemize}\itemsep=0.05cm
      \item El error decae como $O(1/\sqrt{n})$: cuadruplicar $n$ reduce el error máximo a la mitad.
      \item ``$n\ge30$'' es razonable para poblaciones moderadamente simétricas; distribuciones muy asimétricas (con $\rho$ grande) requieren $n$ mayores.
    \end{itemize}
  \end{alertblock}
\end{frame}

% Slide 7: Aplicaciones
\begin{frame}{Aplicaciones: Sumas y Proporciones}
  \footnotesize
  El TLC se aplica igual de bien a sumas y a proporciones muestrales: \pause

  \vspace{0.15cm}
  \begin{itemize}\itemsep=0.12cm
    \item \textbf{Sumas:} si $T=\sum X_i$, entonces $T \approx N(n\mu, n\sigma^2)$ (seguros, control de inventarios). \pause
    \item \textbf{Proporciones:} $\hat{p}=\bar{X}$ para $X_i \sim \text{Bernoulli}(p)$, con $\hat{p} \approx N(p, p(1-p)/n)$ (encuestas, A/B testing). \pause
    \item \textbf{Aproximación Binomial-Normal:} caso particular de sumas de Bernoulli, con corrección de continuidad de Yates. \pause
    \item \textbf{Bootstrap y remuestreo:} el TLC justifica por qué las distribuciones bootstrap se aproximan a la Normal para $n$ grande.
  \end{itemize}
\end{frame}

% Slide 8: Lab Python (Bloque 1)
\begin{frame}[fragile]{Laboratorio en Python: Convergencia desde una Población Asimétrica}
  \scriptsize
  Prueba de Kolmogorov-Smirnov de $Z_n$ contra $N(0,1)$ para una población Exponencial (\texttt{05.02\_central\_limit\_theorem.py}):
  {\lstset{basicstyle=\fontsize{4pt}{4.8pt}\selectfont\ttfamily,aboveskip=0.1em,belowskip=0.1em}
  \lstinputlisting[firstline=17, lastline=35]{../../code/05_distribuciones_muestreo/05.02_central_limit_theorem.py}}
\end{frame}

% Slide 9: Lab Python (Bloque 2)
\begin{frame}[fragile]{Laboratorio en Python: Tasa de Convergencia de Berry-Esseen}
  \scriptsize
  Estimación empírica del error máximo $\sup_z|F_{Z_n}(z)-\Phi(z)|$ y verificación de la tasa $O(1/\sqrt{n})$:
  {\lstset{basicstyle=\fontsize{4pt}{4.8pt}\selectfont\ttfamily,aboveskip=0.1em,belowskip=0.1em}
  \lstinputlisting[firstline=38, lastline=64]{../../code/05_distribuciones_muestreo/05.02_central_limit_theorem.py}}
\end{frame}

% Slide 10: Lab Python (Bloque 3)
\begin{frame}[fragile]{Laboratorio en Python: TLC para Sumas y Proporciones}
  \scriptsize
  Aplicación del TLC a un total de reclamaciones de seguros y a una proporción muestral, con verificación Monte Carlo:
  {\lstset{basicstyle=\fontsize{4pt}{4.8pt}\selectfont\ttfamily,aboveskip=0.1em,belowskip=0.1em}
  \lstinputlisting[firstline=67, lastline=96]{../../code/05_distribuciones_muestreo/05.02_central_limit_theorem.py}}
\end{frame}

% Slide 11: Salida Terminal
\begin{frame}[fragile]{Laboratorio en Python: Salida en Terminal}
  \scriptsize
  Salida estándar generada al ejecutar el script del laboratorio en Python:
  \vspace{0.02cm}
  \begin{block}{Salida Estándar en Consola}
    \fontsize{4pt}{5pt}\selectfont
    \begin{verbatim}
=== Block 1: CLT Convergence from a Skewed Population ===
Exp(lambda=0.25): n=5 D=0.0612 | n=30 D=0.0253 | n=100 D=0.0141 (D shrinks with n)

=== Block 2: Berry-Esseen Rate of Convergence ===
n=10: max gap=0.0418 | n=40: 0.0220 | n=160: 0.0120
Ratio n=10->40: 0.526 (theo 0.500) | n=40->160: 0.545 (theo 0.500)

=== Block 3: CLT for Sums and Proportions ===
Insurance: E(T)=80000, SD(T)=3000, Z=1.667, P(T>85000)=0.0478
Proportion: SD(phat)=0.0324, Z=1.543, P(phat>0.35)=0.0614 (MC: 0.0541)
    \end{verbatim}
  \end{block}
\end{frame}

% Slide 12A: Ejercicio Nivel 1 (Enunciado)
\begin{frame}{Ejercicio en Clase --- Nivel 1: Fundamental (1/2)}
  \footnotesize
  \begin{block}{Problema 5.2.1 --- Aplicación Directa del TLC (Enunciado)}
    Una población tiene media $\mu=200$ y desviación estándar $\sigma=50$. Se toma una muestra aleatoria simple de $n=100$. Use el TLC para aproximar $P(180 < \bar{X} < 210)$.
  \end{block}

  \vspace{0.15cm}
  \pause
  \begin{alertblock}{Planteamiento}
    \begin{itemize}\itemsep=0.04cm
      \item $\text{DE}(\bar{X}) = \sigma/\sqrt{n} = 50/10 = 5$.
      \item Estandarice ambos extremos del intervalo.
    \end{itemize}
  \end{alertblock}
\end{frame}

% Slide 12B: Ejercicio Nivel 1 (Resolución)
\begin{frame}{Ejercicio en Clase --- Nivel 1: Fundamental (2/2)}
  \scriptsize
  Estandarizamos ambos extremos: \pause

  \vspace{0.05cm}
  \begin{block}{Cálculo}
    \begin{align*}
      P(180 < \bar{X} < 210) \approx P\left(\dfrac{180-200}{5} < Z < \dfrac{210-200}{5}\right) = P(-4<Z<2) = \Phi(2)-\Phi(-4) \approx 0.9772.
    \end{align*}
  \end{block}

  \vspace{0.05cm}
  \pause
  \begin{alertblock}{Interpretación}
    \scriptsize
    Como $\Phi(-4)\approx 0$, prácticamente toda la probabilidad relevante viene del extremo superior: hay un $97.7\%$ de probabilidad de que la media muestral esté en el rango indicado.
  \end{alertblock}
\end{frame}

% Slide 13A: Ejercicio Nivel 2 (Enunciado)
\begin{frame}{Ejercicio en Clase --- Nivel 2: Operativo (1/2)}
  \footnotesize
  \begin{block}{Problema 5.2.4 --- TLC para una Suma (Enunciado)}
    Una aseguradora recibe $n=100$ reclamaciones independientes, cada una con media $\mu=\$800$ y desviación estándar $\sigma=\$300$. Use el TLC para aproximar la probabilidad de que el monto total $T=\sum X_i$ exceda $\$85{,}000$.
  \end{block}

  \vspace{0.15cm}
  \pause
  \begin{alertblock}{Estrategia}
    \scriptsize
    Para una suma (no un promedio), use $E(T)=n\mu$ y $\Var(T)=n\sigma^2$ (no dividir entre $n$).
  \end{alertblock}
\end{frame}

% Slide 13B: Ejercicio Nivel 2 (Resolución)
\begin{frame}{Ejercicio en Clase --- Nivel 2: Operativo (2/2)}
  \scriptsize
  Calculamos media y desviación estándar de la suma: \pause

  \vspace{0.05cm}
  \begin{block}{Cálculo}
    \begin{align*}
      E(T) = 100\cdot 800 = 80{,}000, \qquad \text{DE}(T)=\sqrt{100}\cdot 300 = 3000.
    \end{align*}
    \begin{align*}
      P(T>85{,}000) \approx P\left(Z>\dfrac{85{,}000-80{,}000}{3000}\right) = P(Z>1.667) \approx 0.0478.
    \end{align*}
  \end{block}

  \vspace{0.05cm}
  \pause
  \begin{alertblock}{Interpretación}
    \scriptsize
    Hay aproximadamente un $4.78\%$ de probabilidad de que las reclamaciones totales excedan el presupuesto de $\$85{,}000$; la aseguradora puede usar esto para dimensionar reservas de capital.
  \end{alertblock}
\end{frame}

% Slide 14A: Ejercicio Nivel 3 (Enunciado)
\begin{frame}{Ejercicio en Clase --- Nivel 3: Analítico (1/2)}
  \footnotesize
  \begin{block}{Problema 5.2.7 --- TLC para la Exponencial vía FGM (Enunciado)}
    Verifique la demostración del TLC vía FGM específicamente para $X_i \sim \text{Exp}(\lambda)$ (con $\mu=1/\lambda$, $\sigma=1/\lambda$), usando $M_{X_1}(t)=(1-t/\lambda)^{-1}$.
  \end{block}

  \vspace{0.1cm}
  \pause
  \begin{alertblock}{Estrategia}
    \scriptsize
    Sustituya $\mu=\sigma=1/\lambda$ en la fórmula general de $M_{Z_n}(t)$; observe que $\lambda\sigma=1$ simplifica considerablemente la expresión.
  \end{alertblock}
\end{frame}

% Slide 14B: Ejercicio Nivel 3 (Resolución)
\begin{frame}{Ejercicio en Clase --- Nivel 3: Analítico (2/2)}
  \scriptsize
  Sustituimos y simplificamos usando $\lambda\sigma=1$: \pause

  \vspace{0.05cm}
  \begin{block}{Demostración}
    \scriptsize
    \begin{align*}
      M_{Z_n}(t) = e^{-t\sqrt{n}}\left(1-\dfrac{t}{\sqrt{n}}\right)^{-n}.
    \end{align*}
    Con $-\log(1-t/\sqrt n) = t/\sqrt n + t^2/(2n)+O(n^{-3/2})$:
    \begin{align*}
      \log M_{Z_n}(t) = -t\sqrt{n} + n\left[\dfrac{t}{\sqrt{n}}+\dfrac{t^2}{2n}+O(n^{-3/2})\right] = \dfrac{t^2}{2}+O(n^{-1/2}) \to \dfrac{t^2}{2}. \quad \qedhere
    \end{align*}
  \end{block}

  \vspace{0.05cm}
  \pause
  \begin{alertblock}{Interpretación}
    \scriptsize
    Aunque la Exponencial es fuertemente asimétrica, el TLC igual converge: la demostración general no exige simetría, solo media y varianza finitas.
  \end{alertblock}
\end{frame}

% Slide 15A: Ejercicio Nivel 4 (Enunciado)
\begin{frame}{Ejercicio en Clase --- Nivel 4: Desafiante (1/2)}
  \footnotesize
  \begin{block}{Problema 5.2.10 --- TLC con Corrección por Población Finita (Enunciado)}
    Una población finita de $N=500$ cuentas tiene $\sigma=\$80$. Se extrae una muestra de $n=50$ cuentas \emph{sin reemplazo}. Combine el TLC con la FPC (Sección 05.01) para aproximar $P(\bar{X} > \mu + 15)$.
  \end{block}

  \vspace{0.15cm}
  \pause
  \begin{alertblock}{Estrategia}
    \scriptsize
    Primero calcule $\Var(\bar X)$ con FPC, luego aplique el TLC con esa varianza ajustada en lugar de $\sigma^2/n$.
  \end{alertblock}
\end{frame}

% Slide 15B: Ejercicio Nivel 4 (Resolución)
\begin{frame}{Ejercicio en Clase --- Nivel 4: Desafiante (2/2)}
  \scriptsize
  Combinamos FPC (05.01) con el TLC de hoy: \pause

  \vspace{0.05cm}
  \begin{block}{Cálculo}
    \scriptsize
    \begin{align*}
      \Var(\bar{X}) = \dfrac{80^2}{50}\cdot\dfrac{450}{499} = 128\cdot 0.9018 \approx 115.43, \qquad \text{DE}(\bar{X})\approx 10.744.
    \end{align*}
    \begin{align*}
      P(\bar{X}>\mu+15) \approx P\left(Z>\dfrac{15}{10.744}\right) = P(Z>1.396) \approx 0.0814.
    \end{align*}
  \end{block}

  \vspace{0.05cm}
  \pause
  \begin{alertblock}{Interpretación}
    \scriptsize
    Este problema conecta dos secciones: la FPC ajusta la varianza por el muestreo sin reemplazo de una población finita, y el TLC nos permite usar la Normal para calcular la probabilidad final, incluso sin conocer la distribución exacta de las cuentas.
  \end{alertblock}
\end{frame}

% Slide 16: Cierre
\begin{frame}{Cierre de Sección: TLC Asintótico}
  \begin{reflexion}
    El Teorema del Límite Central es, en un sentido preciso, el resultado más importante de la estadística clásica: garantiza que la media muestral converge a la Normal sin importar la forma de la población original. Hoy formalizamos ese resultado con una demostración rigurosa vía FGM y cuantificamos su velocidad de convergencia con el teorema de Berry-Esseen.
  \end{reflexion}

  \vspace{0.2cm}
  \pause
  \begin{block}{Lecciones Clave}
    \begin{itemize}\itemsep=0.08cm
      \item \textbf{TLC:} $Z_n = (\bar{X}_n-\mu)/(\sigma/\sqrt{n}) \xrightarrow{d} N(0,1)$, sin importar la distribución original.
      \item \textbf{Demostración vía FGM:} $\log M_{Z_n}(t) \to t^2/2$ para cualquier población con momentos finitos.
      \item \textbf{Berry-Esseen:} el error de la aproximación decae como $O(1/\sqrt{n})$, cuantificando la heurística ``$n\ge30$''.
    \end{itemize}
  \end{block}
\end{frame}

% Slide 17: Perspectiva Modular
\begin{frame}{Perspectiva Modular --- El Próximo Reto}
  \scriptsize
  \begin{retoclase}
    El TLC aproxima la distribución de $\bar{X}$ usando la Normal cuando $\sigma$ es \emph{conocida}. En la práctica, casi siempre debemos \emph{estimar} $\sigma$ con $S$, lo cual introduce variabilidad adicional. La Sección 05.03 estudiará la distribución Chi-cuadrada, clave para caracterizar la distribución exacta de $S^2$; las Secciones 05.04 y 05.05 introducirán las distribuciones $t$ y $F$, que sustituyen exactamente a la Normal cuando $\sigma$ es desconocida.
  \end{retoclase}

  \vspace{0.08cm}
  \pause
  \begin{alertblock}{Preparación Recomendada}
    \begin{itemize}\itemsep=0.06cm
      \item Repasa la conexión Gamma--Chi-cuadrada (Sección 04.07): $\chi^2_\nu = \text{Gamma}(\nu/2, 2)$.
      \item Reflexiona sobre por qué reemplazar $\sigma$ por $S$ cambia la distribución exacta de $(\bar X-\mu)/(S/\sqrt n)$ de Normal a $t$ de Student.
    \end{itemize}
  \end{alertblock}
\end{frame}

\end{document}
